\section{Exclusion Rules}
\label{sec:rules}

\subsection{Two tiers}

\begin{center}
\begin{tabular}{@{}p{0.30\linewidth}p{0.30\linewidth}p{0.30\linewidth}@{}}
\toprule
 & \emph{ignore} (hide) & \emph{deny} (block) \\
\midrule
Purpose & Reduce clutter & Never serve \\
File & \code{\~{}/.config/fsb/ignore} & \code{\~{}/.config/fsb/deny} \\
Listings and search & Entry omitted & Entry omitted \\
Direct access (read, download, preview, details) & Allowed & Blocked, answered as 404 \\
Negation (\code{!}) & Allowed & Refused \\
Defaults & \code{.DS\_Store}, \code{.git/}, \code{node\_modules/} & The core rules (Section~\ref{sec:core}) \\
\bottomrule
\end{tabular}
\end{center}

\begin{req}{RUL-1}
A path matched by a deny rule is not listed, not found by search, not readable,
downloadable, previewable or describable, and not reachable through any other
endpoint.  A request for it is answered exactly as a request for a path that does
not exist: the same status and the same body, except under \code{--debug}
(\rid{CLI-6}).  Deny always wins over ignore; no ignore rule, including a negation,
can expose a denied path.
\end{req}

\begin{req}{RUL-2}
A path matched by an ignore rule is omitted from listings and search results but
remains reachable by its path.  Entering an ignored folder on purpose (by typing its
path or following a link) lists its contents; only the folder's own entry, in its
parent's listing, is hidden.
\end{req}

\begin{req}{RUL-3}
A denied folder denies everything inside it.  Once a folder is excluded, no later rule
can bring anything inside it back.
\end{req}

\begin{req}{RUL-4}
Symbolic links and aliases do not change what a rule means.  A path is refused if
\emph{either} the path as requested \emph{or} the real location it resolves to is
denied; a link whose target is denied is refused, and is not listed.  The children of
a folder are judged by the folder's real location, however it was reached.  On macOS
the aliases of the data volume are the same as the ordinary paths:
\code{/System/Volumes/Data/Users/me} is \code{/Users/me}, in any letter case and
any spelling that the file system equates (\rid{NAM-1}).
\end{req}

\begin{req}{RUL-5}
Every spelling that the file system resolves to a denied file is refused.  This
includes other letter cases, decomposed accents, characters that fold to ASCII
(\rid{NAM-1}), doubled or trailing separators, dot segments, symbolic-link and
firmlink aliases, and resource-fork names such as \code{file/..namedfork/rsrc}.
\end{req}

\begin{req}{RUL-6}
A file with more than one name (a hard link) is judged by its identity: if any of
its names is denied, it is refused under every name and is not listed.  The other
names are found under the roots and under the credential locations of every home
folder on the machine, using an index built in the background at startup (about 20
seconds for a home folder of 700{,}000 files).  Until the index is ready, and for a
file with several names that it has not yet seen, the file is refused rather than
guessed.  A denied name that lies outside every root and every credential location
is not detected.
\end{req}

\begin{req}{RUL-7}
A permission error on a path that lies inside a denied place is reported as not
found, so that an unreadable file in a denied folder cannot be told from a missing
one.  A permission error elsewhere is reported as such (status 403) and stays visible
to the user.
\end{req}

\begin{req}{RUL-8}
Denying is decided before anything is read.  A denied path never causes a cloud-only
file to be downloaded.
\end{req}

\subsection{Rule syntax}
\label{sec:rulefiles}

The syntax is that of \code{.gitignore}, with the additions and restrictions below.

\begin{req}{RUL-9}
A line is blank, a comment (first character \code{\#}), or a rule.  Surrounding white
space is ignored.  A leading byte order mark and any line ending (LF, CR LF or bare
CR) are accepted.  In a rule, \code{*} matches any characters except \code{/},
\code{?} one character except \code{/}, \code{[abc]} and \code{[!abc]} a class
(which never matches \code{/}), and \code{**} any number of folders.  A backslash
makes the next character literal (\code{\textbackslash\#} for a literal \code{\#}).  A
trailing \code{/} makes the rule apply to folders only.
\end{req}

\begin{req}{RUL-10}
A pattern with no \code{/} matches a name at any depth.  A pattern that begins with
\code{\~{}/} is relative to the home folder; one that begins with \code{/} is
absolute; one that contains a \code{/} elsewhere is relative to the home folder; and
one that begins with \code{**/} matches at any depth.
\end{req}

\begin{req}{RUL-11}
Matching ignores letter case and normalization exactly as in \rid{NAM-1}, and ignores
trailing white space in a name (macOS permits names such as \code{id.pem\textvisiblespace}
and \code{.env} followed by a tab), and \code{**} matches names that contain newlines.
\end{req}

\begin{req}{RUL-12}
A rule file must work as written or fail loudly.  \texttt{fsb} refuses to start,
naming the file, the line and the problem, for: negation in the deny file; an
empty pattern; a trailing backslash; a \code{/} inside \code{[...]}; \code{\~{}name}
(another user's home; only \code{\~{}} and \code{\~{}/...} are supported); a
pattern wrapped in quotes; an inline comment (a \code{\#} after white space; write
\code{\textbackslash\#} for a literal one); an invisible or formatting character; and a
character that looks like a slash but is not one.
\end{req}

\subsection{Core and optional deny rules}
\label{sec:core}

\begin{req}{RUL-13}
The core deny rules cover locations and files that hold credentials or secrets.  They
are active by default and built in, so they apply even when no deny file exists.
\code{fsb --init} writes them, uncommented, into the deny file so that the user can
see them.  They match at any depth, so the same folders in a backup or clone of the
home folder, in another user's home, or under a wrong \code{HOME} are protected too.
\end{req}

\begin{req}{RUL-14}
The core rules are:

\begin{center}
\begin{tabular}{@{}p{0.45\linewidth}p{0.45\linewidth}@{}}
\toprule
Folders & Files \\
\midrule
\code{.ssh/}, \code{.aws/}, \code{.gnupg/}, \code{.kube/} &
\code{.netrc}, \code{.git-credentials}, \code{.pgpass} \\
\code{**/.config/gh/} &
\code{.env}, \code{.env.*} \\
\code{**/Library/Keychains/} &
\code{*.pem}, \code{*.key} \\
\code{**/Library/Application Support/Google/Chrome/} &
\code{id\_rsa}, \code{id\_dsa}, \code{id\_ecdsa}, \code{id\_ed25519}, \code{id\_ecdsa\_sk}, \code{id\_ed25519\_sk} \\
\code{**/Library/Application Support/Firefox/} &
\code{*.p12}, \code{*.pfx}, \code{*.ppk}, \code{*.jks}, \code{*.keystore} \\
\code{**/Library/Safari/}, \code{**/Library/Cookies/} &
\code{*.kdbx}, \code{*.keychain}, \code{*.keychain-db} \\
\bottomrule
\end{tabular}
\end{center}

Public key files (\code{id\_rsa.pub} and the like) are not denied.  The list is
versioned with the release, and additions are called out in the release notes.
\end{req}

\begin{req}{RUL-15}
The secret-file patterns are deliberately broad and have known costs.  \code{*.key}
also matches Keynote presentations; \code{.env} also matches a folder named
\code{.env} (such as a Python virtual environment) and everything in it; and
\code{.env.*} also matches committed templates such as \code{.env.example}.  A user who
needs any of them removes the pattern from the deny file and accepts the warning of
\rid{RUL-16}.
\end{req}

\begin{req}{RUL-16}
When a deny file exists but does not contain one or more core rules (they were removed
or commented out, or a release added a rule the file predates), \texttt{fsb}
(a) prints the warning of \rid{CLI-11}, for example
\code{fsb: warning: core deny rule(s) disabled: .ssh/, .aws/; these paths are
browsable.}, and (b) shows a persistent banner in the interface with the same
message.  Both disappear once every core rule is in effect.  Rules are compared as
names are (\rid{NAM-1}), so \code{.SSH/} counts as the core rule \code{.ssh/}, but a
rule that differs in meaning does not.  The condition cannot be suppressed by any flag.
\end{req}

\begin{req}{RUL-17}
The optional rules, shipped commented out for the user to enable, are situational or
prone to false positives: \code{\~{}/.docker/config.json}, \code{\~{}/.npmrc},
\code{\~{}/Library/Mail/} and \code{\~{}/Library/Messages/}.
\end{req}
